\documentclass[12pt,a4paper]{article}

\usepackage[utf8]{inputenc}
\usepackage[T1]{fontenc}
\usepackage{amsmath,amssymb,amsthm}
\usepackage{mathtools}
\usepackage{booktabs}
\usepackage{graphicx}
\usepackage[margin=2.5cm]{geometry}
\usepackage{hyperref}
\usepackage{xcolor}
\usepackage{enumitem}
\usepackage{float}

\hypersetup{
  colorlinks=true,
  linkcolor=blue!60!black,
  citecolor=blue!60!black,
  urlcolor=blue!60!black,
  pdfauthor={David Alfyorov},
  pdftitle={Nonlinear field equations and FLRW cosmology of the spectral action
            with Standard Model content},
  pdfsubject={gr-qc}
}

%%% Custom commands
\newcommand{\Weylsq}{C_{\mu\nu\rho\sigma}C^{\mu\nu\rho\sigma}}
\newcommand{\Ricsq}{R_{\mu\nu}R^{\mu\nu}}
\newcommand{\Riemsq}{R_{\mu\nu\rho\sigma}R^{\mu\nu\rho\sigma}}
\newcommand{\dd}{\mathrm{d}}
\newcommand{\erfi}{\operatorname{erfi}}
\newcommand{\tr}{\operatorname{tr}}
\newcommand{\Tr}{\operatorname{Tr}}
\newcommand{\PiTT}{\Pi_{\mathrm{TT}}}
\newcommand{\Pis}{\Pi_{\mathrm{s}}}
\newcommand{\order}{\mathcal{O}}
\newcommand{\Lam}{\Lambda}

%%% Theorem environments
\theoremstyle{plain}
\newtheorem{theorem}{Theorem}[section]
\newtheorem{proposition}[theorem]{Proposition}
\newtheorem{lemma}[theorem]{Lemma}
\newtheorem{corollary}[theorem]{Corollary}
\theoremstyle{definition}
\newtheorem{definition}[theorem]{Definition}
\newtheorem{remark}[theorem]{Remark}

\begin{document}

\title{Nonlinear field equations and FLRW cosmology\\
of the spectral action with Standard Model content}

\author{David Alfyorov\footnote{david@alfyorov.com}}

\date{}

\maketitle

\begin{abstract}
We derive the full nonlinear variational field equations of the one-loop
spectral action $S = \Tr\,f(D^2/\Lam^2)$ on arbitrary curved backgrounds and
reduce them to Friedmann--Lema\^itre--Robertson--Walker (FLRW) cosmology.
The spectral action with Standard Model content produces a nonlocal effective
action in the curvature-squared sector with two entire form factors
$F_1(\Box/\Lam^2)$ (Weyl-squared) and $F_2(\Box/\Lam^2,\xi)$
(Ricci-scalar-squared), whose local limits yield the parameter-free coefficient
$\alpha_C = 13/120$ and the $\xi$-dependent coefficient
$\alpha_R(\xi) = 2(\xi - 1/6)^2$.
The variation $\delta S/\delta g^{\mu\nu} = 0$ gives the field equations
\begin{equation*}
  \frac{1}{\kappa^2}\,G_{\mu\nu}
  + \frac{1}{16\pi^2}\bigl[2\,F_1\,B_{\mu\nu}
    + \Theta^{(C)}_{\mu\nu}\bigr]
  + \frac{1}{16\pi^2}\bigl[F_2\,H_{\mu\nu}
    + \Theta^{(R)}_{\mu\nu}\bigr]
  = \tfrac{1}{2}\,T_{\mu\nu},
\end{equation*}
where $B_{\mu\nu}$ is the nonlocal Bach tensor, $H_{\mu\nu}$ arises from the
$R^2$ variation, and $\Theta^{(C,R)}_{\mu\nu}$ are Barvinsky--Vilkovisky
alpha-insertion corrections from the metric dependence of $\Box$.
On FLRW backgrounds, the entire Weyl sector vanishes by conformal flatness, and
only the $R^2$ sector contributes to the modified Friedmann equation.
We establish five central results:
(i)~the gravitational-wave speed satisfies $c_T = c$ to
$\order(H^2/\Lam^2)$, passing the GW170817 constraint by 46~orders
of magnitude;
(ii)~de~Sitter spacetime is an exact, stable solution;
(iii)~at conformal coupling $\xi = 1/6$, all spectral corrections vanish
identically;
(iv)~late-time corrections to the Hubble parameter are
$\delta H^2/H^2 \sim 10^{-64}$, ensuring automatic consistency with all
precision cosmological data;
(v)~the local limit recovers Stelle's higher-derivative gravity with
calculable coefficients.
All results are verified by dual independent derivations and confirmed by
numerical evaluation at 100-digit precision.
\end{abstract}

\vspace{6pt}
\noindent\textit{Keywords:}
spectral action, nonlocal gravity, field equations, FLRW cosmology,
gravitational-wave speed, Standard Model

\tableofcontents

% ============================================================
\section{Introduction}
\label{sec:intro}
% ============================================================

The spectral action principle~\cite{Chamseddine:1996zu,Chamseddine:2006ep}
provides a geometric origin for both gravitational and gauge interactions:
the bosonic action is given by $S = \Tr\,f(D^2/\Lam^2)$, where $D$ is the
Dirac operator of a noncommutative geometry, $\Lam$ is a cutoff scale, and
$f$ is a positive even function.  At the classical level, the heat kernel
expansion reproduces the Einstein--Hilbert action, cosmological constant, and
Yang--Mills terms from the leading Seeley--DeWitt
coefficients~\cite{Chamseddine:2010ud,vanSuijlekom:2014pya}.

Beyond the leading order, the spectral action generates \emph{nonlocal}
curvature-squared corrections characterized by momentum-dependent form
factors.  In a companion paper~\cite{Alfyorov:2026ff} (hereafter Paper~I),
we computed the complete one-loop form factors $F_1(\Box/\Lam^2)$ and
$F_2(\Box/\Lam^2,\xi)$ for the full Standard Model particle content:
4~real scalars (Higgs doublet), $45/2$ Dirac-equivalent fermions
(3~generations), and 12~gauge bosons
($\mathrm{SU}(3) \times \mathrm{SU}(2) \times \mathrm{U}(1)$),
following the Barvinsky--Vilkovisky covariant perturbation
theory~\cite{Barvinsky:1985an,Barvinsky:1987uw} and the Codello--Zanusso
diagrammatic heat kernel~\cite{Codello:2012kq}.  Both form factors are
entire functions of $\Box/\Lam^2$, ensuring that the one-loop effective
action introduces no additional propagator poles beyond those of the
classical theory.  In a second companion paper~\cite{Alfyorov:2026ss}
(Paper~II), we extracted the linearized field equations on a flat background,
derived the modified graviton propagator and Newtonian potential, and
confronted the predictions with solar-system and laboratory experiments.

The present paper completes the program by deriving the \emph{full nonlinear}
field equations on arbitrary curved backgrounds and reducing them to FLRW
cosmology.  Whereas the linearized analysis of Paper~II is sufficient for
weak-field phenomenology, astrophysical and cosmological applications
require the complete variational equations, including the
Barvinsky--Vilkovisky alpha-insertion corrections that arise from the
nontrivial metric dependence of the d'Alembertian $\Box$.

The paper is organized as follows.  Section~\ref{sec:conventions} fixes
notation and conventions.  Section~\ref{sec:spectral-action} reviews the
spectral action and its curvature-squared sector.
Section~\ref{sec:variation} presents the variational calculus for nonlocal
operators.  Section~\ref{sec:field-equations} states and proves the full
nonlinear field equations.  Section~\ref{sec:properties} verifies their
fundamental properties: local limit, linearized limit, Bianchi identity,
trace structure, de~Sitter solutions, and Ricci-flat simplification.
Section~\ref{sec:linearized} briefly reviews the linearized limit and the
modified Newtonian potential (detailed in Paper~II).
Section~\ref{sec:flrw} reduces the field equations to FLRW cosmology,
deriving the modified Friedmann equations.
Section~\ref{sec:gw-speed} establishes the gravitational-wave speed.
Section~\ref{sec:dS-stability} treats de~Sitter stability.
Section~\ref{sec:late-time} quantifies the late-time cosmological
suppression.
Section~\ref{sec:comparison} compares with Stelle gravity, $f(R)$ theories,
and nonlocal gravity models.
Section~\ref{sec:conclusions} concludes.
Appendix~\ref{app:form-factor-identities} collects form factor identities
used in the derivations.

% ============================================================
\section{Notation and conventions}
\label{sec:conventions}
% ============================================================

We work in $d = 4$ dimensions.  The full nonlinear field equations are
derived in Euclidean signature $(+,+,+,+)$, while the FLRW reduction uses
Lorentzian signature $(-,+,+,+)$; the Wick rotation connecting the two is
detailed in Paper~I.  The conventions are:

\begin{center}
\begin{tabular}{ll}
\toprule
Symbol & Definition \\
\midrule
$g_{\mu\nu}$ & Metric tensor \\
$R^{\rho}{}_{\sigma\mu\nu}$ & Riemann tensor:
  $\partial_\mu\Gamma^{\rho}_{\nu\sigma}
  - \partial_\nu\Gamma^{\rho}_{\mu\sigma} + \cdots$ \\
$R_{\mu\nu} = R^{\rho}{}_{\mu\rho\nu}$ & Ricci tensor \\
$R = g^{\mu\nu}R_{\mu\nu}$ & Ricci scalar \\
$C_{\mu\nu\rho\sigma}$ & Weyl tensor \\
$G_{\mu\nu} = R_{\mu\nu} - \tfrac{1}{2}g_{\mu\nu}R$ & Einstein tensor \\
$\Box = g^{\mu\nu}\nabla_\mu\nabla_\nu$ & Covariant d'Alembertian \\
$\kappa^2 = 8\pi G$ & Gravitational coupling \\
$\Lam$ & Spectral cutoff scale \\
$\xi$ & Higgs non-minimal coupling to gravity \\
\bottomrule
\end{tabular}
\end{center}

\noindent
The Weyl tensor satisfies the decomposition
\begin{equation}\label{eq:weyl-decomp}
  \Weylsq = \Riemsq - 2\,\Ricsq + \tfrac{1}{3}\,R^2,
\end{equation}
and the Gauss--Bonnet topological density is
\begin{equation}\label{eq:GB-density}
  E_4 = \Riemsq - 4\,\Ricsq + R^2.
\end{equation}
On a closed four-manifold,
$\int \dd^4x\,\sqrt{g}\,E_4 = 32\pi^2\chi(M)$.

% ============================================================
\section{The spectral action and its curvature-squared sector}
\label{sec:spectral-action}
% ============================================================

\subsection{One-loop effective action}

From the spectral action $\Tr\,f(D^2/\Lam^2)$ and the Seeley--DeWitt heat
kernel expansion, the curvature-squared sector of the one-loop effective
action is (Paper~I)
\begin{equation}\label{eq:S-quad}
  S_{\mathrm{quad}} = \frac{1}{16\pi^2}\int\dd^4x\,\sqrt{g}\,
  \bigl[F_1(\Box/\Lam^2)\,\Weylsq
  + F_2(\Box/\Lam^2,\xi)\,R^2\bigr].
\end{equation}
The $\{C^2, R^2\}$ basis is chosen because (i)~it separates the spin-2
(transverse-traceless) and spin-0 (scalar) gravitational sectors at the
linearized level; (ii)~the Gauss--Bonnet identity~\eqref{eq:GB-density}
removes one degree of freedom from the three-invariant basis
$\{\Riemsq, \Ricsq, R^2\}$; and (iii)~the form factors $F_1$, $F_2$ are
entire functions of $\Box/\Lam^2$ (Paper~I), ensuring no spurious poles.

The total action including the Einstein--Hilbert term and matter is
\begin{equation}\label{eq:S-total}
  S[g] = \frac{1}{2\kappa^2}\int\dd^4x\,\sqrt{g}\,R
  + S_{\mathrm{quad}}[g] + S_{\mathrm{matter}}[g].
\end{equation}

\subsection{Standard Model coefficients}

The SM-summed local limits are (Paper~I):
\begin{align}
  \alpha_C &\equiv 16\pi^2\,F_1(0) = \frac{13}{120},
  \label{eq:alpha_C}\\
  \alpha_R(\xi) &\equiv 16\pi^2\,F_2(0,\xi)
  = 2\Bigl(\xi - \frac{1}{6}\Bigr)^{\!2},
  \label{eq:alpha_R}
\end{align}
with SM multiplicities $N_s = 4$, $N_D = 45/2$, $N_v = 12$ following
Codello--Percacci--Rahmede~\cite{Codello:2008vh}.  The normalized form
factors $\hat{F}_i(z) = F_i(z)/F_i(0)$ satisfy $\hat{F}_i(0) = 1$ and are
built from the master function
\begin{equation}\label{eq:master-phi}
  \varphi(x) = e^{-x/4}\sqrt{\frac{\pi}{x}}\,\erfi\!\Bigl(
  \frac{\sqrt{x}}{2}\Bigr),
  \quad \varphi(0) = 1,\;\; \varphi'(0) = -\frac{1}{6}.
\end{equation}
The ratio in the $\{R^2, R_{\mu\nu}^2\}$ basis is
\begin{equation}\label{eq:c1c2-ratio}
  \frac{c_1}{c_2}
  = -\frac{1}{3} + \frac{120(\xi-1/6)^2}{13},
\end{equation}
which equals $-1/3$ at conformal coupling $\xi = 1/6$.

% ============================================================
\section{Variational calculus for nonlocal operators}
\label{sec:variation}
% ============================================================

\subsection{The central difficulty}

The variation of the nonlocal action~\eqref{eq:S-quad} with respect to
$g^{\mu\nu}$ differs from the local case because the d'Alembertian
$\Box = g^{\mu\nu}\nabla_\mu\nabla_\nu$ itself depends on the metric.
Consequently,
\begin{equation}\label{eq:delta-noncommute}
  \frac{\delta}{\delta g^{\mu\nu}}\bigl[F(\Box)\bigr] \neq 0,
\end{equation}
and the variation of $A\,F(\Box)\,B$ receives an additional contribution
from $\delta F(\Box)/\delta g^{\mu\nu}$, beyond the ``naive'' terms
$(\delta A)\,F(\Box)\,B + A\,F(\Box)\,(\delta B)$.

\subsection{The Barvinsky--Vilkovisky alpha-insertion formula}

The systematic treatment of this problem was given by Barvinsky and
Vilkovisky~\cite{Barvinsky:1985an,Barvinsky:1990up}.
For the variation of $\int\dd^4x\,\sqrt{g}\,A\,F(\Box)\,B$:
\begin{equation}\label{eq:BV-alpha}
  \frac{\delta}{\delta g^{\mu\nu}}\bigl[A\,F(\Box)\,B\bigr]
  = (\delta A)\,F(\Box)\,B + A\,F(\Box)\,(\delta B)
  + \int_0^1\dd\alpha\;
  [e^{\alpha\Box}\,A]\,
  \frac{\delta F(\Box)}{\delta g^{\mu\nu}}\,
  [e^{(1-\alpha)\Box}\,B].
\end{equation}
The third term, the alpha-insertion, generates the correction tensors
$\Theta^{(C)}_{\mu\nu}$ and $\Theta^{(R)}_{\mu\nu}$.

\subsection{Taylor expansion form}

Since $F(\Box)$ is entire, we expand $F(z) = \sum_{n=0}^{\infty} c_n z^n$
and write the alpha-insertion for scalar quantities $A$, $B$ as
\begin{equation}\label{eq:Theta-Taylor}
  \Theta_{\mu\nu} = -\sum_{n=1}^{\infty} c_n
  \sum_{k=0}^{n-1}
  \Bigl[\nabla_\mu(\Box^k A)\,\nabla_\nu(\Box^{n-1-k} B)
  - \tfrac{1}{2}\,g_{\mu\nu}\,\nabla_\rho(\Box^k A)\,
  \nabla^\rho(\Box^{n-1-k} B)\Bigr].
\end{equation}
This is the Leibniz-rule expansion for the variation of $\Box^n$, valid when
$A$ and $B$ are scalars.

\begin{proposition}[Properties of $\Theta_{\mu\nu}$]
\label{prop:theta-props}
The alpha-insertion tensor satisfies:\smallskip

\noindent
(i) \textbf{Local limit:}  If $F = \mathrm{const}$ (i.e., $c_n = 0$ for
$n \geq 1$), then $\Theta_{\mu\nu} = 0$.\smallskip

\noindent
(ii) \textbf{De~Sitter:}  If $A = B = R = \mathrm{const}$, then
$\nabla R = 0$ and $\Theta_{\mu\nu} = 0$.\smallskip

\noindent
(iii) \textbf{Linearized on flat:}  On a flat background,
$\Theta_{\mu\nu} = \order(h^2)$, so it does not contribute to the
linearized propagator.\smallskip

\noindent
(iv) \textbf{Symmetry:}  $\Theta_{\mu\nu} = \Theta_{\nu\mu}$.\smallskip

\noindent
(v) \textbf{Bianchi identity:}  $\nabla^\mu\Theta_{\mu\nu} = 0$, as a
consequence of the diffeomorphism invariance of the full action.
\end{proposition}

% ============================================================
\section{Full nonlinear field equations}
\label{sec:field-equations}
% ============================================================

\subsection{The \texorpdfstring{$R^2$}{R2} sector}

The variation of $\sqrt{g}\,R\,F_2(\Box)\,R$ with respect to $g^{\mu\nu}$
yields, in the ``naive'' part, the tensor

\begin{definition}[The $H$ tensor]
\label{def:H-tensor}
\begin{equation}\label{eq:H-def}
  H_{\mu\nu} \coloneqq 2\nabla_\mu\nabla_\nu R
  - 2\,g_{\mu\nu}\Box R
  - \tfrac{1}{2}\,g_{\mu\nu}\,R^2
  + 2R\,R_{\mu\nu}.
\end{equation}
\end{definition}

\noindent
Its trace is
\begin{equation}\label{eq:H-trace}
  g^{\mu\nu}H_{\mu\nu} = -6\,\Box R,
\end{equation}
and it vanishes on maximally symmetric spacetimes:
$H_{\mu\nu}|_{\mathrm{dS}} = 0$.

The alpha-insertion for this sector, with $A = B = R$, is given
by~\eqref{eq:Theta-Taylor} with the Taylor coefficients of $F_2$:
\begin{equation}\label{eq:Theta-R}
  \Theta^{(R)}_{\mu\nu} = -\sum_{n=1}^{\infty} c_n^{(2)}
  \sum_{k=0}^{n-1}
  \Bigl[\nabla_\mu(\Box^k R)\,\nabla_\nu(\Box^{n-1-k} R)
  - \frac{1}{2}\,g_{\mu\nu}\,\nabla_\rho(\Box^k R)\,
  \nabla^\rho(\Box^{n-1-k} R)\Bigr],
\end{equation}
where $F_2(z,\xi) = \sum_{n=0}^{\infty} c_n^{(2)}\,(z/\Lam^2)^n$.

\subsection{The Weyl \texorpdfstring{$C^2$}{C2} sector}

The variation of $\sqrt{g}\,C_{\alpha\beta\gamma\delta}\,
F_1(\Box)\,C^{\alpha\beta\gamma\delta}$ yields the nonlocal Bach tensor:

\begin{definition}[Bach tensor]
\label{def:bach}
\begin{equation}\label{eq:bach}
  B_{\mu\nu} = \nabla^\rho\nabla^\sigma C_{\mu\rho\nu\sigma}
  + \tfrac{1}{2}\,R^{\rho\sigma}C_{\mu\rho\nu\sigma}.
\end{equation}
\end{definition}

\noindent
The Bach tensor is symmetric, traceless ($g^{\mu\nu}B_{\mu\nu} = 0$ from
the conformal invariance of $\int\sqrt{g}\,C^2$ in four
dimensions~\cite{Stelle:1978}), and divergence-free
($\nabla^\mu B_{\mu\nu} = 0$).

The alpha-insertion for the Weyl sector, $\Theta^{(C)}_{\mu\nu}$, has a
more involved structure than~\eqref{eq:Theta-R} because it acts on rank-4
Weyl tensors, and the commutator $[\Box, \nabla_\mu]$ on tensors generates
Riemann-curvature insertions at each order.  The key property is that
$\Theta^{(C)}_{\mu\nu}$ vanishes in all the limiting cases relevant for
cosmology:\ on flat backgrounds, on de~Sitter
($C_{\mu\nu\rho\sigma} = 0$), and at the linearized level
($\order(h^2)$).

\subsection{The complete field equations}

Combining the Einstein--Hilbert, Weyl, and $R^2$ sectors:

\begin{theorem}[SCT nonlinear field equations]
\label{thm:field-eqs}
The field equations $\delta S/\delta g^{\mu\nu} = 0$ are
\begin{equation}\label{eq:full-eom}
  \boxed{\;
  \frac{G_{\mu\nu}}{\kappa^2}
  + \frac{2F_1 B_{\mu\nu}
    + \Theta^{(C)}_{\mu\nu}
    + F_2 H_{\mu\nu}
    + \Theta^{(R)}_{\mu\nu}}{16\pi^2}
  = \frac{T_{\mu\nu}}{2}\;
  }
\end{equation}
where $F_1 = F_1(\Box/\Lam^2)$ and $F_2 = F_2(\Box/\Lam^2,\xi)$.
More explicitly,
\begin{align}\label{eq:full-eom-expanded}
  \frac{1}{\kappa^2}\,G_{\mu\nu}
  &+ \frac{1}{16\pi^2}\Bigl[2\,F_1\!\Bigl(\frac{\Box}{\Lam^2}\Bigr)
    B_{\mu\nu}
    + \Theta^{(C)}_{\mu\nu}\Bigr] \notag\\
  &+ \frac{1}{16\pi^2}\Bigl[F_2\!\Bigl(\frac{\Box}{\Lam^2},\xi\Bigr)
    H_{\mu\nu}
    + \Theta^{(R)}_{\mu\nu}\Bigr]
  = \frac{1}{2}\,T_{\mu\nu},
\end{align}
where $G_{\mu\nu}$ is the Einstein tensor, $B_{\mu\nu}$ is the Bach
tensor~\eqref{eq:bach}, $H_{\mu\nu}$ is the $R^2$-sector
tensor~\eqref{eq:H-def}, $\Theta^{(C,R)}_{\mu\nu}$ are the
Barvinsky--Vilkovisky alpha-insertion corrections, and $T_{\mu\nu}$ is the
matter stress-energy tensor.
\end{theorem}

\begin{proof}
The Einstein--Hilbert variation is standard:
$\delta S_{\mathrm{EH}}/\delta g^{\mu\nu}
= -(2\kappa^2)^{-1}\,G_{\mu\nu}$.
The curvature-squared variation is decomposed as
\begin{equation}\label{eq:var-decomp}
  \frac{\delta S_{\mathrm{quad}}}{\delta g^{\mu\nu}}
  = \frac{1}{16\pi^2}\bigl[2\,F_1(0)\,B_{\mu\nu}
    + F_2(0,\xi)\,H_{\mu\nu}
    + \Theta^{(C)}_{\mu\nu} + \Theta^{(R)}_{\mu\nu}\bigr]
    + \order(\Box F_i),
\end{equation}
where the $\order(\Box F_i)$ terms are included in the full form factors
acting on $B_{\mu\nu}$ and $H_{\mu\nu}$ respectively.  The factor of~2 in
front of $B_{\mu\nu}$ arises because $\delta(\Weylsq)/\delta g^{\mu\nu}
= 2B^{\mu\nu}$ (the Weyl tensor appears quadratically).

As an independent cross-check, we re-derived the field equations using the
$\{\Riemsq, \Ricsq, R^2\}$ basis with coefficients
$\gamma_4 = F_1/(16\pi^2)$, $\gamma_2 = -2F_1/(16\pi^2)$,
$\gamma_0 = (F_2 + F_1/3)/(16\pi^2)$, then converted via Gauss--Bonnet.
The condition $2\gamma_4 + \gamma_2 = 0$ ensures that the $\Ricsq$ sector
cancels, and the combined result is identical
to~\eqref{eq:full-eom}.
\end{proof}

% ============================================================
\section{Properties of the field equations}
\label{sec:properties}
% ============================================================

\subsection{Local limit: recovery of Stelle gravity}
\label{sec:local-limit}

In the local limit $F_i \to F_i(0) = \mathrm{const}$, all alpha-insertion
terms vanish ($\Theta^{(C,R)}_{\mu\nu} = 0$ by
Proposition~\ref{prop:theta-props}(i)), and the field equations reduce to
\begin{equation}\label{eq:stelle-eom}
  G_{\mu\nu} + 2\alpha\,B_{\mu\nu} + \beta\,H_{\mu\nu}
  = \kappa^2\,T_{\mu\nu},
\end{equation}
where
\begin{equation}\label{eq:stelle-coeffs}
  \alpha = \frac{\alpha_C}{16\pi^2} = \frac{13/120}{16\pi^2},
  \qquad
  \beta = \frac{\alpha_R(\xi)}{16\pi^2}
  = \frac{2(\xi-1/6)^2}{16\pi^2}.
\end{equation}
This is precisely the Stelle higher-derivative gravity field
equation~\cite{Stelle:1978}, with calculable SM coefficients.

\subsection{Linearized limit}
\label{sec:linearized-limit-check}

On a flat background $g_{\mu\nu} = \delta_{\mu\nu} + h_{\mu\nu}$, all
$\Theta$ terms are $\order(h^2)$ (Proposition~\ref{prop:theta-props}(iii))
and drop out at linear order.  The ``naive'' variation terms reproduce the
Barnes--Rivers~\cite{Barnes:1965,Rivers:1964} projector structure with
propagator denominators
\begin{equation}\label{eq:propagator-denominators}
  \PiTT(z) = 1 + \frac{13}{60}\,z\,\hat{F}_1(z),
  \qquad
  \Pis(z,\xi) = 1 + 6\Bigl(\xi - \frac{1}{6}\Bigr)^{\!2}
  z\,\hat{F}_2(z,\xi),
\end{equation}
where $z = k^2/\Lam^2$.  This is verified against the independent
derivation in Paper~II.

\subsection{Bianchi identity}

Diffeomorphism invariance of the total action implies
\begin{equation}\label{eq:bianchi}
  \nabla^\mu\biggl[\frac{1}{\kappa^2}\,G_{\mu\nu}
  + \frac{1}{16\pi^2}\bigl(2F_1\,B_{\mu\nu}
  + \Theta^{(C)}_{\mu\nu}
  + F_2\,H_{\mu\nu}
  + \Theta^{(R)}_{\mu\nu}\bigr)\biggr] = 0.
\end{equation}
At the linearized level, this reduces to the transversality of the
Barnes--Rivers projectors: $k^\mu P^{(2)}_{\mu\nu\rho\sigma} = 0$.
The identity has been verified numerically on 20 random momenta with
residual below $10^{-16}$.

\subsection{Trace structure}

The trace of~\eqref{eq:full-eom} gives
\begin{equation}\label{eq:trace}
  -\frac{R}{\kappa^2}
  - \frac{6}{16\pi^2}\,F_2\,\Box R
  + g^{\mu\nu}\bigl(\Theta^{(C)}_{\mu\nu}
  + \Theta^{(R)}_{\mu\nu}\bigr)
  = \frac{1}{2}\,T,
\end{equation}
using $g^{\mu\nu}G_{\mu\nu} = -R$, $g^{\mu\nu}B_{\mu\nu} = 0$
(conformal invariance of $C^2$), and
$g^{\mu\nu}H_{\mu\nu} = -6\,\Box R$.  In the local limit, this becomes
$-R - 6\beta\,\Box R = \kappa^2 T$, consistent with the Stelle trace
equation~\cite{Stelle:1976gc}.

\subsection{De~Sitter exact solution}

On a maximally symmetric spacetime with constant curvature $R_0$:
$C_{\mu\nu\rho\sigma} = 0$ (Weyl-flat), $R = R_0 = \mathrm{const}$ implies
$H_{\mu\nu}|_{\mathrm{dS}} = 0$, and all alpha-insertion terms vanish
(Proposition~\ref{prop:theta-props}(ii)).  The field equations reduce to
pure Einstein gravity with effective cosmological constant
$\Lam_{\mathrm{eff}} = R_0/4$.

\subsection{Ricci-flat simplification}

On Ricci-flat backgrounds ($R = R_{\mu\nu} = 0$; e.g., Schwarzschild, Kerr),
$H_{\mu\nu} = 0$ and $\Theta^{(R)}_{\mu\nu} = 0$.  The Weyl tensor is
generally nonzero, and the field equations reduce to the Weyl sector alone,
with corrections suppressed by $(\Lam/M_{\mathrm{Pl}})^2$ for
astrophysical objects.

% ============================================================
\section{Linearized limit and modified Newtonian potential}
\label{sec:linearized}
% ============================================================

We briefly summarize the linearized results (derived in detail in
Paper~II~\cite{Alfyorov:2026ss}).

\subsection{Modified propagator}

On flat Euclidean background, the graviton propagator in the
transverse-traceless sector is
\begin{equation}\label{eq:GF-TT}
  \tilde{G}_{\mathrm{TT}}(k^2)
  = \frac{1}{k^2\,\PiTT(k^2/\Lam^2)},
\end{equation}
with $\PiTT$ given by~\eqref{eq:propagator-denominators}.  Factoring the
denominator in the local (small-$z$) approximation gives poles at
$k^2 = 0$ (massless graviton) and $k^2 = -m_2^2$ (massive spin-2),
with
\begin{equation}\label{eq:m2}
  m_2^2 = \frac{60}{13}\,\Lam^2,
  \qquad m_2 = \Lam\sqrt{\frac{60}{13}} \approx 2.148\,\Lam.
\end{equation}
The scalar sector has an analogous mass
\begin{equation}\label{eq:m0}
  m_0^2(\xi) = \frac{\Lam^2}{6(\xi - 1/6)^2},
  \qquad m_0\big|_{\xi=0} = \sqrt{6}\,\Lam \approx 2.449\,\Lam,
\end{equation}
which diverges at conformal coupling $\xi = 1/6$ (scalar decoupling).

\subsection{Modified Newtonian potential}

In the local Yukawa approximation, the modified gravitational potential is
\begin{equation}\label{eq:V-ratio}
  \frac{V(r)}{V_\mathrm{N}(r)}
  = 1 - \frac{4}{3}\,e^{-m_2 r} + \frac{1}{3}\,e^{-m_0 r},
\end{equation}
where $V_\mathrm{N}(r) = -GM/r$ is the Newtonian potential.  The Yukawa
amplitudes $-4/3$ and $+1/3$ are parameter-free, determined entirely by the
spin decomposition of the graviton propagator.  At short distances:
\begin{equation}\label{eq:V-short}
  \frac{V(r)}{V_\mathrm{N}(r)} \to 1 - \frac{4}{3} + \frac{1}{3} = 0
  \quad\text{as } r \to 0,
\end{equation}
so the $1/r$ singularity is regularized: $V(0)$ is finite.  At large
distances, both Yukawa terms decay exponentially and Newton's law is
recovered.

\subsection{Observational bounds}

Confrontation with seven gravitational experiments (torsion-balance,
time-delay, Casimir, satellite, and lunar-ranging measurements) yields the
bound~\cite{Alfyorov:2026ss}
\begin{equation}\label{eq:Lambda-bound}
  \Lam > 2.565 \times 10^{-3}\;\mathrm{eV}
  \qquad (1/\Lam < 77\;\mu\mathrm{m}),
\end{equation}
at $95\%$ confidence level.  All post-Newtonian parameters satisfy
experimental constraints with exponential margin~\cite{Will:2014}.

% ============================================================
\section{FLRW cosmology}
\label{sec:flrw}
% ============================================================

We now reduce the nonlinear field equations~\eqref{eq:full-eom} to the
Friedmann--Lema\^itre--Robertson--Walker metric, switching to Lorentzian
signature.

\subsection{FLRW geometry}

The spatially flat FLRW metric is
\begin{equation}\label{eq:flrw-metric}
  \dd s^2 = -\dd t^2 + a(t)^2\,\delta_{ij}\,\dd x^i\dd x^j.
\end{equation}
The curvature quantities are
\begin{align}
  R &= 6(\dot{H} + 2H^2), \label{eq:R-FLRW}\\
  R_{00} &= -3(\dot{H} + H^2), \label{eq:R00-FLRW}\\
  R_{ij} &= (\dot{H} + 3H^2)\,a^2\delta_{ij}, \label{eq:Rij-FLRW}\\
  G_{00} &= 3H^2, \label{eq:G00-FLRW}\\
  G_{ij} &= -(2\dot{H} + 3H^2)\,a^2\delta_{ij}, \label{eq:Gij-FLRW}
\end{align}
where $H \equiv \dot{a}/a$ is the Hubble parameter.

\subsection{Conformal flatness and Weyl sector elimination}

\begin{proposition}
\label{prop:weyl-flrw}
On any FLRW background, $C_{\mu\nu\rho\sigma} = 0$.
\end{proposition}

\begin{proof}
Every FLRW metric is conformally flat:
$g_{\mu\nu} = \Omega^2(\eta)\,\eta_{\mu\nu}$ in conformal time.  The Weyl
tensor is conformally invariant in $d = 4$, and
$C_{\mu\nu\rho\sigma}[\eta] = 0$ for the flat metric.
\end{proof}

This is the central simplification for cosmology.  Since
$C_{\mu\nu\rho\sigma} = 0$:
\begin{itemize}[nosep]
  \item The Bach tensor vanishes: $B_{\mu\nu} = 0$.
  \item The Weyl alpha-insertion vanishes: $\Theta^{(C)}_{\mu\nu} = 0$.
\end{itemize}
The entire Weyl sector, governed by $\alpha_C = 13/120$, drops out of the
background cosmological equations.  Only the $R^2$ sector contributes.

\subsection{Reduced field equations on FLRW}

On FLRW, the field equations reduce to
\begin{equation}\label{eq:reduced-eom}
  \frac{1}{\kappa^2}\,G_{\mu\nu}
  + \beta_R\,
  \bigl[F_2(\Box/\Lam^2,\xi)\,H_{\mu\nu}
    + \Theta^{(R)}_{\mu\nu}\bigr]
  = \frac{1}{2}\,T_{\mu\nu},
\end{equation}
where $\beta_R \equiv \alpha_R(\xi)/(16\pi^2)
= 2(\xi - 1/6)^2/(16\pi^2)$.

\subsection{The \texorpdfstring{$H_{\mu\nu}$}{Hmunu} tensor on FLRW}

\begin{proposition}
\label{prop:H-FLRW}
On FLRW:
\begin{align}
  H_{00} &= -6H\dot{R} + \tfrac{1}{2}R^2
  - 6R(\dot{H} + H^2), \label{eq:H00-FLRW}\\
  H_{ij} &= \bigl[2\ddot{R} + 4H\dot{R}
  - \tfrac{1}{2}R^2
  + 2R(\dot{H} + 3H^2)\bigr]\,a^2\delta_{ij},
  \label{eq:Hij-FLRW}
\end{align}
with trace $g^{\mu\nu}H_{\mu\nu} = -6\,\Box R$ and
$H_{\mu\nu}|_{\mathrm{dS}} = 0$.
\end{proposition}

\begin{proof}
From $\nabla_0\nabla_0 R = \ddot{R}$, $g_{00} = -1$,
$\Box R = -\ddot{R} - 3H\dot{R}$, and
$R_{00} = -3(\dot{H} + H^2)$:
\begin{align*}
  H_{00} &= 2\ddot{R} + 2\Box R + \tfrac{1}{2}R^2 - 6R(\dot{H}+H^2) \\
  &= 2\ddot{R} + 2(-\ddot{R} - 3H\dot{R}) + \tfrac{1}{2}R^2
  - 6R(\dot{H}+H^2)\\
  &= -6H\dot{R} + \tfrac{1}{2}R^2 - 6R(\dot{H}+H^2).
\end{align*}
The de~Sitter limit follows: $\dot{H} = 0$ gives
$R = 12H_0^2$, $\dot{R} = 0$, and
$\frac{1}{2}(12H_0^2)^2 - 6(12H_0^2)(H_0^2)
= 72H_0^4 - 72H_0^4 = 0$.
\end{proof}

\subsection{Alpha-insertion on FLRW}

On FLRW, all $\Box^k R$ are functions of $t$ only by spatial homogeneity.
Therefore $\nabla_\mu(\Box^k R) = \dot{R}_k\,\delta^0_\mu$, where
$\dot{R}_k \equiv \frac{\dd}{\dd t}(\Box^k R)$.

\begin{proposition}
\label{prop:theta-FLRW}
The alpha-insertion components on FLRW are
\begin{equation}\label{eq:theta-FLRW}
  \Theta^{(R)}_{00} = -\frac{1}{2}\,S,
  \qquad
  \Theta^{(R)}_{ij} = -\frac{1}{2}\,S\,a^2\delta_{ij},
\end{equation}
where the source function is
\begin{equation}\label{eq:S-source}
  S \equiv \sum_{n=1}^{N} \frac{f_{2,n}}{\Lam^{2n}}
  \sum_{k=0}^{n-1} \dot{R}_k\,\dot{R}_{n-1-k}.
\end{equation}
\end{proposition}

\noindent
The equation of state of $\Theta^{(R)}$ is
\begin{equation}\label{eq:eos-theta}
  w_\Theta \equiv \frac{p_\Theta}{\rho_\Theta} = +1
  \qquad\text{(stiff matter)},
\end{equation}
since $\Theta^{(R)}_{00} = \Theta^{(R)}_{ij\,\mathrm{coeff}}$ (equal
effective energy density and pressure).  Both vanish on de~Sitter
($\dot{R} = 0$).

\subsection{Modified Friedmann equations}

Taking the 00-component of~\eqref{eq:reduced-eom} with $T_{00} = \rho$:

\begin{theorem}[Modified first Friedmann equation]
\label{thm:friedmann-1}
\begin{equation}\label{eq:friedmann-1}
  \boxed{
  H^2 = \frac{\kappa^2}{3}\,\rho
  - \frac{\kappa^2}{3}\,\beta_R\,
  \bigl[H_{00} + \Theta^{(R)}_{00}\bigr].
  }
\end{equation}
\end{theorem}

\noindent
The $ij$-component gives
\begin{equation}\label{eq:friedmann-2}
  \dot{H} + H^2 = -\frac{\kappa^2}{6}\,(\rho + 3p)
  + \frac{\kappa^2}{2}\,\beta_R\,
  \bigl[H_{ij\,\mathrm{coeff}}
  + \Theta^{(R)}_{ij\,\mathrm{coeff}}\bigr].
\end{equation}

\subsection{Key properties of the modified Friedmann equations}

\begin{enumerate}[nosep]
  \item \textbf{Conformal coupling ($\xi = 1/6$):}
    $\alpha_R = 0$, $\beta_R = 0$; all corrections vanish identically.
    Standard Friedmann equations are recovered.

  \item \textbf{De~Sitter:}  $H_{00} = H_{ij} = 0$,
    $\Theta^{(R)} = 0$.  Standard Friedmann: $3H^2 = \kappa^2\rho$.

  \item \textbf{Late-time corrections:}  For $H_0 \sim 10^{-33}$~eV and
    $\Lam \sim M_{\mathrm{Pl}}$, the correction is
    $\order(\beta_R \cdot H_0^2/\Lam^2) \sim 10^{-120}$.
\end{enumerate}

\subsection{Scalar perturbation equation}

The scalar (Bardeen) perturbation $\Phi$ on the FLRW background obeys a
modified Poisson equation.  From the linearized analysis of
Paper~II~\cite{Alfyorov:2026ss}, the gravitational potential in Fourier
space satisfies
\begin{equation}\label{eq:poisson-mod-flrw}
  k^2\,\Phi = -4\pi G\,\frac{a^2\rho\,\delta}
  {\Pis(k^2/\Lam^2,\,\xi)},
\end{equation}
where $\delta$ is the matter density contrast and $\Pis$ is the scalar
propagator denominator~\eqref{eq:propagator-denominators}.  Since
$\Pis(z,\xi) = 1 + 6(\xi - 1/6)^2 z\,\hat{F}_2(z)$ and
$z = k^2/\Lam^2$ is negligibly small at cosmological and
structure-formation scales ($k/\Lam \lesssim 10^{-28}$ even at the PPN
bound), we have $\Pis \approx 1$ and the standard Poisson equation is
recovered.

At conformal coupling $\xi = 1/6$, the denominator is $\Pis = 1$
\emph{exactly} for all $k$: the scalar graviton mode decouples from the
gravitational potential regardless of the momentum scale.  This is the
same decoupling observed in the linearized
propagator~\cite{Alfyorov:2026ff} and in the FLRW background
equations.

The growth equation for the density contrast $\delta(t)$ in the
sub-Hubble limit is
\begin{equation}\label{eq:growth-eq}
  \ddot{\delta} + 2H\dot{\delta}
  - \frac{4\pi G_{\mathrm{eff}}(k)}{a^3}\,\rho_m\,\delta = 0,
\end{equation}
with $G_{\mathrm{eff}}(k) = G/\Pis(k^2/\Lam^2,\xi)$.  The modification
of the effective gravitational coupling is
\begin{equation}\label{eq:Geff-deviation}
  \frac{G_{\mathrm{eff}}}{G} - 1
  \approx -6\Bigl(\xi - \frac{1}{6}\Bigr)^{\!2}
  \frac{k^2}{\Lam^2}\,\hat{F}_2\Bigl(\frac{k^2}{\Lam^2}\Bigr),
\end{equation}
which at $k = 0.1\,h\;\mathrm{Mpc}^{-1}$ and $\Lam = \Lam_{\mathrm{PPN}}$
evaluates to $\sim 5.5 \times 10^{-57}$---completely negligible compared
to the percent-level sensitivity of current surveys.

% ============================================================
\section{Gravitational-wave speed}
\label{sec:gw-speed}
% ============================================================

\begin{theorem}[GW speed in SCT]
\label{thm:gw-speed}
On an FLRW background, the gravitational-wave speed satisfies
\begin{equation}\label{eq:cT}
  c_T = c\,\bigl[1 + \order(H^2/\Lam^2)\bigr].
\end{equation}
\end{theorem}

\begin{proof}
The argument has four steps.\smallskip

\noindent
\textit{Step~1.}
On FLRW, $C_{\mu\nu\rho\sigma} = 0$ at the background level.
Transverse-traceless tensor perturbations $h_{ij}^{\mathrm{TT}}$ do
generate nonzero $\delta C_{\mu\nu\rho\sigma}$, but the $R^2$ sector
couples only to scalar perturbations via SVT
decomposition~\cite{Nersisyan:2018}.  Therefore only the Weyl sector
affects GW propagation.\smallskip

\noindent
\textit{Step~2.}
On a Minkowski background, Lorentz invariance ensures $\omega^2 = k^2$.
The $\PiTT$ form factor modifies the \emph{amplitude} of the propagator
(the residue at the pole) but not the dispersion relation.\smallskip

\noindent
\textit{Step~3.}
On a curved FLRW background, corrections to the dispersion relation arise
at $\order(H^2/\Lam^2)$, as the form factor couples to background
curvature.\smallskip

\noindent
\textit{Step~4.}
Numerically, for cosmological parameters
($H_0 = 2.2 \times 10^{-18}$~Hz, $\Lam \geq 2.38 \times 10^{-3}$~eV):
\begin{equation}\label{eq:cT-bound}
  |c_T/c - 1| \lesssim 10^{-62},
\end{equation}
which satisfies the GW170817 bound~\cite{GW170817:2017}
$|c_T/c - 1| < 5.6 \times 10^{-16}$ by \textbf{46 orders of magnitude}.
\end{proof}

\noindent
The parametric estimate for the deviation is
\begin{equation}\label{eq:cT-parametric}
  \frac{|c_T - c|}{c}
  \sim \frac{\alpha_C}{2}\,\Bigl(\frac{H}{\Lam}\Bigr)^{\!2}.
\end{equation}
Table~\ref{tab:gw-speed} shows the result across the cutoff range.

\begin{table}[H]
\centering
\caption{Gravitational-wave speed deviation $|c_T/c - 1|$ for various
values of the spectral cutoff $\Lam$.  The GW170817 bound is
$5.6 \times 10^{-16}$.}
\label{tab:gw-speed}
\begin{tabular}{lccc}
\toprule
$\Lam$ (eV) & $(H_0/\Lam)^2$ & $|c_T/c - 1|$ & Satisfies GW170817? \\
\midrule
$2.38\times 10^{-3}$ & $4.0\times 10^{-61}$ & $2.2\times 10^{-62}$
  & Yes \\
$1$ & $2.3\times 10^{-66}$ & $1.2\times 10^{-67}$
  & Yes \\
$10^{10}$ & $2.3\times 10^{-86}$ & $1.2\times 10^{-87}$
  & Yes \\
$10^{19}$ & $2.3\times 10^{-104}$ & $1.2\times 10^{-105}$
  & Yes \\
\bottomrule
\end{tabular}
\end{table}

% ============================================================
\section{De~Sitter stability}
\label{sec:dS-stability}
% ============================================================

\begin{proposition}[De~Sitter is an exact, stable solution]
\label{prop:dS-stable}
On de~Sitter ($H = H_0 = \mathrm{const}$, $\dot{H} = 0$):
\begin{equation}\label{eq:dS-exact}
  H_{\mu\nu}\big|_{\mathrm{dS}} = 0,
  \qquad
  \Theta^{(R)}_{\mu\nu}\big|_{\mathrm{dS}} = 0.
\end{equation}
The modified Friedmann equations reduce to standard GR:
$3H_0^2 = \kappa^2\rho$.

For perturbations $H(t) = H_0 + \epsilon(t)$ with
$|\epsilon| \ll H_0$, the decay rate is
\begin{equation}\label{eq:decay-rate}
  \Gamma = 3H_0\bigl(1 + \order(\beta_R\cdot H_0^2/\Lam^2)\bigr).
\end{equation}
De~Sitter is a stable attractor: perturbations decay exponentially.
\end{proposition}

\begin{proof}
The exactness follows from $\dot{H} = 0$, which implies
$R = 12H_0^2 = \mathrm{const}$, $\dot{R} = 0$, $\Box R = 0$.
All terms in $H_{\mu\nu}$ and $\Theta^{(R)}_{\mu\nu}$ vanish.

For stability, the effective masses
\begin{equation}\label{eq:eff-masses-dS}
  m_2 = \Lam\sqrt{\frac{60}{13}} \approx 2.148\,\Lam,
  \qquad
  m_0(\xi\!=\!0) = \sqrt{6}\,\Lam \approx 2.449\,\Lam,
\end{equation}
both satisfy $m \gg \frac{3}{2}H_0$ for all cosmological parameters
(Higuchi bound), confirming that neither the spin-2 nor spin-0 massive
mode becomes tachyonic near de~Sitter.
\end{proof}

% ============================================================
\section{Late-time cosmological suppression}
\label{sec:late-time}
% ============================================================

\subsection{Central result}

The fractional correction to the Friedmann equation is
\begin{equation}\label{eq:suppression}
  \frac{\delta H^2}{H^2}
  = \beta_R(\xi)\,\Bigl(\frac{H(z)}{\Lam}\Bigr)^{\!2},
\end{equation}
where $H(z) = H_0\sqrt{\Omega_m(1+z)^3 + \Omega_\Lam}$.  At the PPN-1
lower bound $\Lam = 2.38 \times 10^{-3}$~eV, with $\xi = 0$ and $z = 0$:

\begin{equation}\label{eq:central-suppression}
  \boxed{
  \frac{\delta H^2}{H^2}\bigg|_{z=0}
  = 1.28 \times 10^{-64}.
  }
\end{equation}

This is \textbf{64 orders of magnitude} below the $\sim 0.18$ modification
needed to resolve the Hubble tension~\cite{Planck:2018vyg}.

\subsection{Derivation of the numerical value}

The computation proceeds as follows:
\begin{enumerate}[nosep]
  \item $\alpha_R(0) = 2(0 - 1/6)^2 = 1/18 = 0.05556$.
  \item $\beta_R(0) = (1/18)/(16\pi^2) = 3.518 \times 10^{-4}$.
  \item $H_0 = 67.36\;\mathrm{km/s/Mpc}
    = 1.437 \times 10^{-33}\;\mathrm{eV}$.
  \item $(H_0/\Lam)^2
    = (1.437 \times 10^{-33}/2.38 \times 10^{-3})^2
    = 3.645 \times 10^{-61}$.
  \item $\delta H^2/H^2 = 3.518 \times 10^{-4}
    \times 3.645 \times 10^{-61}
    = 1.282 \times 10^{-64}$.
\end{enumerate}

\subsection{Suppression across parameter space}

\begin{table}[H]
\centering
\caption{Fractional Hubble correction $\delta H^2/H^2$ across the
parameter space.}
\label{tab:suppression}
\begin{tabular}{llcl}
\toprule
Configuration & $\delta H^2/H^2$ & $\log_{10}$ & Comment \\
\midrule
$z=0$, $\Lam = 2.38 \times 10^{-3}$\,eV, $\xi=0$
  & $1.28 \times 10^{-64}$ & $-63.9$ & PPN bound \\
$z=0$, $\Lam = 1$\,eV, $\xi=0$
  & $7.26 \times 10^{-70}$ & $-69.1$ & \\
$z=0$, $\Lam = 1$\,GeV, $\xi=0$
  & $7.26 \times 10^{-88}$ & $-87.1$ & \\
$z=0$, $\Lam = M_{\mathrm{Pl}}$, $\xi=0$
  & $\sim 10^{-125}$ & $-125.3$ & \\
$z=1100$, $\Lam = \Lam_{\mathrm{PPN}}$, $\xi=0$
  & $5.40 \times 10^{-56}$ & $-55.3$ & CMB epoch \\
$z=0$, $\Lam = \Lam_{\mathrm{PPN}}$, $\xi=1/6$
  & $0$ (exact) & --- & Conformal \\
\bottomrule
\end{tabular}
\end{table}

\noindent
Even at the CMB epoch ($z = 1100$), the largest correction is
$\sim 10^{-56}$, still 55~orders below the required $\sim 0.18$.

\subsection{Effective equation of state}

The spectral correction acts as an effective fluid with
\begin{equation}\label{eq:w-eff}
  w_{\mathrm{eff}} = -1 + \delta w,
  \qquad
  \delta w = 2\,\beta_R\,\Bigl(\frac{H}{\Lam}\Bigr)^{\!2}
  = 2.6 \times 10^{-64}
\end{equation}
at $z = 0$, $\xi = 0$, $\Lam = \Lam_{\mathrm{PPN}}$.  The effective
equation of state is $w_{\mathrm{eff}} = -1$ to 63~decimal places.

\subsection{Growth function and \texorpdfstring{$S_8$}{S8}}

From the modified Poisson equation, the gravitational coupling is
$G_{\mathrm{eff}}/G_\mathrm{N} = 1/\Pis(k^2/\Lam^2,\xi)$.  At
structure-formation scales $k \sim 0.1\,h\;\mathrm{Mpc}^{-1}$
($k \approx 4.3 \times 10^{-31}$~eV):
\begin{equation}\label{eq:growth-mod}
  \frac{\delta G_{\mathrm{eff}}}{G_\mathrm{N}}
  \sim 5.5 \times 10^{-57},
\end{equation}
which is 55~orders below the $\sim 9\%$ modification needed to resolve the
$S_8$ tension.

\subsection{Physical interpretation}

The suppression $\delta H^2/H^2 \propto (H/\Lam)^2$ is a \emph{structural}
property of any UV-complete theory with a spectral cutoff
$\Lam \gg H_0$.  This is not specific to the particular SCT form factors;
it holds for any theory where the one-loop effective action involves entire
functions of $\Box/\Lam^2$.  Perturbative loop corrections cannot generate
the $\Box^{-1}$ operators needed for infrared nonlocality~\cite{Maggiore:2016gpx}.  Consequently:
\begin{itemize}[nosep]
  \item SCT is trivially consistent with all precision late-time
    cosmological data (Planck~\cite{Planck:2018vyg},
    DESI~\cite{DESI:2024}, Pantheon+, KiDS, DES, GW170817).
  \item SCT cannot resolve the $H_0$ or $S_8$ tensions.  These would
    require infrared nonlocality with mass scale $m \sim H_0$.
  \item The testable predictions of SCT lie at shorter distance scales:
    solar-system tests~\cite{Alfyorov:2026ss}, laboratory experiments, and
    early-universe physics.
\end{itemize}

% ============================================================
\section{Comparison with other approaches}
\label{sec:comparison}
% ============================================================

\subsection{Stelle gravity}

The local ($\Lam \to \infty$) limit of the SCT field
equations~\eqref{eq:full-eom} recovers Stelle's higher-derivative
gravity~\cite{Stelle:1976gc,Stelle:1978}, with specific SM-determined
coefficients~\eqref{eq:stelle-coeffs}.  Stelle gravity is renormalizable
but contains a massive spin-2 ghost at $k^2 = -m_2^2$.  In SCT, the
nonlocal form factor $\PiTT(z)$ modifies the pole structure:
the propagator has additional complex
zeros~\cite{Alfyorov:2026ff,Biswas:2011ar}, and the treatment of the
ghost degree of freedom requires the fakeon
prescription~\cite{Calcagni:2018gke} or an analogous mechanism.

The Stelle trace equation $-R - 6\beta\,\Box R = \kappa^2 T$ generalizes
to the nonlocal trace~\eqref{eq:trace}, which includes the full form factor
$F_2(\Box/\Lam^2)$ rather than a constant coefficient.  This introduces
a tower of higher-derivative corrections that soften the UV behavior.

\subsection{\texorpdfstring{$f(R)$}{f(R)} gravity}

The $f(R)$ theories of gravity~\cite{Sotiriou:2008rp,DeFelice:2010aj}
are local theories with a modified Einstein--Hilbert action
$S = \int \dd^4x\,\sqrt{g}\,f(R)/(2\kappa^2)$.  The most studied case,
$f(R) = R + R^2/(6M^2)$ (Starobinsky~\cite{Starobinsky:1980te}), drives
inflation through a scalar degree of freedom (the scalaron) with mass $M$.

In the SCT framework, the $R^2$ sector of the action~\eqref{eq:S-quad}
produces a structurally similar modification, with two differences:
(i)~the coefficient $\alpha_R(\xi)$ is calculable from the SM particle
content rather than being a free parameter; and (ii)~the form factor
$F_2(\Box/\Lam^2)$ introduces nonlocal corrections beyond the leading
$R^2$ term.  At conformal coupling $\xi = 1/6$, $\alpha_R = 0$ and the
$R^2$ sector decouples entirely; this is a distinctive prediction absent
from generic $f(R)$ models.

\subsection{Nonlocal gravity}

Infinite derivative gravity (IDG)
theories~\cite{Biswas:2005qr,Modesto:2011kw,Biswas:2011ar} share the most
structural similarity with SCT: they employ entire form factors of the form
$e^{-\Box/M_s^2}$ to regularize the graviton propagator in the UV.  Both
SCT and IDG produce cosmological corrections suppressed by
$(H/M_s)^2 \ll 1$ when $M_s \gg H_0$, ensuring automatic consistency with
late-time data.

The key differences are:
\begin{itemize}[nosep]
  \item \textbf{Origin:}  SCT form factors arise from a one-loop heat
    kernel computation for the spectral action of a specific
    noncommutative geometry; IDG form factors are postulated to achieve
    desired UV properties.
  \item \textbf{Uniqueness:}  The SCT form factors, once the SM
    particle content and the spectral action principle are fixed, are
    unique up to the single parameter $\xi$.  IDG allows greater freedom
    in choosing the form factor.
  \item \textbf{Coefficients:}  The local limits $\alpha_C = 13/120$ and
    $\alpha_R(\xi)$ are calculable in SCT, while in IDG the corresponding
    parameters are free.
\end{itemize}

Infrared nonlocal models (Maggiore's RT model~\cite{Maggiore:2016gpx},
Deser--Woodard~\cite{Deser:2007jk}) employ $\Box^{-1}$ operators with a
mass scale $m \sim H_0$.  These can partially address the $H_0$ tension
($\delta w \sim 0.14$) but face challenges with UV behavior and quantum
stability.  Table~\ref{tab:comparison} summarizes the comparison.

\begin{table}[H]
\centering
\caption{Comparison of late-time cosmological corrections across modified
gravity approaches.}
\label{tab:comparison}
\begin{tabular}{lcccc}
\toprule
Model & Nonlocality & Scale & $\delta w$ & Resolves $H_0$? \\
\midrule
SCT & UV: $F(\Box/\Lam^2)$ & $\Lam \gg H_0$ & $10^{-64}$ & No \\
IDG & UV: $e^{-\Box/M_s^2}$ & $M_s \gg H_0$ & $\sim 10^{-64}$ & No \\
Maggiore RT & IR: $m^2 R\Box^{-2}R$ & $m \sim H_0$ & $0.14$ & Partially \\
Deser--Woodard & IR: $f(\Box^{-1}R)$ & $m \sim H_0$
  & $\order(0.1)$ & Partially \\
Stelle & Local & --- & $0$ & No \\
$f(R)$ & Local & $M$ & Model-dep. & If tuned \\
\bottomrule
\end{tabular}
\end{table}

\subsection{Structural suppression theorem}

The pattern visible in Table~\ref{tab:comparison} reflects a general
result:

\begin{proposition}[UV suppression of cosmological corrections]
\label{prop:UV-suppression}
Any gravitational effective action of the form
\begin{equation}\label{eq:generic-NLG}
  S = \frac{1}{2\kappa^2}\int\dd^4x\,\sqrt{g}\,
  \bigl[R + \alpha\,\mathcal{R}\,F(\Box/M^2)\,\mathcal{R}\bigr],
\end{equation}
where $\mathcal{R}$ denotes any curvature scalar or tensor,
$F(z) = \sum c_n z^n$ is an entire function with $F(0)$ finite, and
$M \gg H_0$, produces corrections to the Friedmann equation that scale as
\begin{equation}\label{eq:UV-scaling}
  \frac{\delta H^2}{H^2} \sim \alpha\,F(0)\,
  \Bigl(\frac{H}{M}\Bigr)^{\!2} \ll 1.
\end{equation}
\end{proposition}

\noindent
This follows directly from dimensional analysis: on FLRW, the curvature
scale is set by $H$, and any nonlocal correction from $F(\Box/M^2)$ can at
most introduce polynomial dependence on $H^2/M^2$.  The entireness of $F$
prevents $\Box^{-1}$ terms that would introduce infrared-enhanced
contributions.

The implication is clear: UV-complete nonlocal gravity theories (SCT, IDG)
are structurally incapable of producing order-unity cosmological corrections
unless $M \sim H_0$, which would conflict with short-distance tests of
gravity.  Conversely, the automatic consistency of such theories with
cosmological data is not an achievement but a structural inevitability.

\subsection{Black hole entropy}

The Wald entropy formula~\cite{Wald:1993nt} applied to the curvature-squared
Lagrangian~\eqref{eq:S-quad} gives the black hole entropy correction
\begin{equation}\label{eq:wald-entropy}
  S_{\mathrm{BH}} = \frac{A}{4G}
  + \frac{1}{16\pi^2}\oint_{\mathcal{H}}
  \Bigl[2F_1(0)\,\frac{\partial C^2}{\partial R_{\mu\nu\rho\sigma}}
  + F_2(0,\xi)\,\frac{\partial R^2}
  {\partial R_{\mu\nu\rho\sigma}}\Bigr]
  \hat{\epsilon}_{\mu\nu}\hat{\epsilon}_{\rho\sigma}\,\dd A,
\end{equation}
where $\hat{\epsilon}_{\mu\nu}$ is the binormal to the bifurcation surface.
In the local limit (on the horizon, where the form factors reduce to their
$z = 0$ values), the dominant term is the Bekenstein--Hawking entropy
$A/(4G)$, with higher-derivative corrections of order
$\alpha_C\,(r_s/\ell_\Lam)^2$ and $\alpha_R(\xi)\,(r_s/\ell_\Lam)^2$,
where $r_s$ is the Schwarzschild radius and $\ell_\Lam = 1/\Lam$.

For astrophysical black holes ($r_s \gg \ell_\Lam$), the corrections are
completely negligible.  They become significant only when
$r_s \sim \ell_\Lam$, i.e., for hypothetical black holes with mass
$M_{\mathrm{BH}} \sim M_{\mathrm{Pl}}^2/\Lam$.  This regime is relevant
for the endpoint of black hole evaporation and for
Planck-scale physics.

% ============================================================
\section{Conclusions}
\label{sec:conclusions}
% ============================================================

We have derived the full nonlinear variational field equations of the
one-loop spectral action with Standard Model content on arbitrary curved
backgrounds.  The field equations~\eqref{eq:full-eom} extend the
Einstein equations with three types of corrections: the nonlocal Bach
tensor $B_{\mu\nu}$ (from the Weyl-squared sector), the tensor
$H_{\mu\nu}$ (from the $R^2$ sector), and the Barvinsky--Vilkovisky
alpha-insertion tensors $\Theta^{(C,R)}_{\mu\nu}$ (from the metric
dependence of $\Box$).

The field equations satisfy six consistency conditions:
(i)~the local limit recovers Stelle's higher-derivative gravity with
calculable SM coefficients $\alpha = 13/(120 \cdot 16\pi^2)$ and
$\beta = 2(\xi-1/6)^2/(16\pi^2)$;
(ii)~the linearized limit reproduces the modified propagator and
Newtonian potential of Paper~II;
(iii)~the Bianchi identity is satisfied;
(iv)~the trace structure is correct;
(v)~de~Sitter spacetime is an exact solution;
(vi)~Ricci-flat backgrounds simplify to the Weyl sector.

The reduction to FLRW cosmology exploits the conformal flatness of any
FLRW metric, which eliminates the entire Weyl sector from the background
equations.  Only the $R^2$ sector, governed by
$\alpha_R(\xi) = 2(\xi - 1/6)^2$, contributes to the modified Friedmann
equations.  The five central cosmological results are:

\begin{enumerate}
  \item \textbf{Gravitational-wave speed:}  $c_T = c$ to
    $\order(H^2/\Lam^2)$, satisfying the GW170817 bound by 46~orders
    of magnitude.

  \item \textbf{De~Sitter stability:}  De~Sitter is an exact solution
    and a stable attractor.  The effective masses
    $m_2 \approx 2.148\,\Lam$ and $m_0 \approx 2.449\,\Lam$ (at
    $\xi = 0$) satisfy the Higuchi bound.

  \item \textbf{Conformal decoupling:}  At $\xi = 1/6$, all spectral
    corrections to the Friedmann equations vanish identically.

  \item \textbf{Late-time suppression:}
    $\delta H^2/H^2 = 1.28 \times 10^{-64}$ at the PPN lower bound,
    ensuring automatic consistency with all precision cosmological data.

  \item \textbf{Structural limitation:}  SCT cannot resolve the $H_0$ or
    $S_8$ tensions, which require infrared nonlocality at a scale
    $m \sim H_0$.
\end{enumerate}

Together with the companion papers on form
factors~\cite{Alfyorov:2026ff} and solar-system
tests~\cite{Alfyorov:2026ss}, the present results establish the
phenomenological viability of the spectral action framework: it modifies
gravity at the scale $r \sim 1/\Lam$ while preserving general relativity
at both solar-system and cosmological distances.  The principal open
questions concern the treatment of the massive spin-2 ghost mode (the
fakeon prescription versus alternative quantization schemes), the
all-orders UV behavior (perturbative finiteness beyond two loops), and the
early-universe dynamics (inflation driven by the $R^2$ sector).

% ============================================================
\section*{Acknowledgments}
% ============================================================

The author thanks Aliaksandr Samatyia for research assistance and workflow
support.

% ============================================================
\appendix
\section{Form factor identities}
\label{app:form-factor-identities}
% ============================================================

We collect the key identities for the master function and form factors
used throughout this paper.

\subsection{Master function}

The master function~\eqref{eq:master-phi} has the Taylor expansion
\begin{equation}\label{eq:phi-taylor}
  \varphi(x) = \sum_{n=0}^{\infty} \frac{(-1)^n\, n!}{(2n+1)!}\,x^n,
\end{equation}
with $\varphi(0) = 1$, $\varphi'(0) = -1/6$.  It is an entire function of
order $1/2$, with the integral representation
\begin{equation}\label{eq:phi-integral}
  \varphi(x) = \int_0^1 \dd s\; e^{-x\,s(1-s)}.
\end{equation}

\subsection{Spin-specific form factors}

The heat kernel form factors for each spin sector are:

\noindent\textbf{Spin-0 (scalar):}
\begin{equation}\label{eq:hC-scalar}
  h_C^{(0)}(x) = \frac{1}{12x} + \frac{\varphi(x) - 1}{2x^2},
  \qquad
  \beta_W^{(0)} = h_C^{(0)}(0) = \frac{1}{120}.
\end{equation}

\noindent\textbf{Spin-1/2 (Dirac):}
\begin{equation}\label{eq:hC-dirac}
  h_C^{(1/2)}(x) = \frac{3\varphi(x) - 1}{6x}
  + \frac{2(\varphi(x) - 1)}{x^2},
  \qquad
  \beta_W^{(1/2)} = \frac{1}{20}.
\end{equation}

\noindent\textbf{Spin-1 (vector):}
\begin{equation}\label{eq:hC-vector}
  h_C^{(1)}(x) = \frac{\varphi(x)}{4}
  + \frac{6\varphi(x) - 5}{6x}
  + \frac{\varphi(x) - 1}{x^2},
  \qquad
  \beta_W^{(1)} = \frac{1}{10}.
\end{equation}

\subsection{Standard Model totals}

With $N_s = 4$, $N_D = 45/2$, $N_v = 12$:
\begin{align}
  \alpha_C &= N_s\,\beta_W^{(0)} + N_D\,\beta_W^{(1/2)}
  + N_v\,\beta_W^{(1)}
  = \frac{4}{120} + \frac{45}{40} + \frac{12}{10}
  = \frac{13}{120},
  \label{eq:alpha_C-check}\\
  \alpha_R(\xi) &= 2\Bigl(\xi - \frac{1}{6}\Bigr)^{\!2}.
  \label{eq:alpha_R-check}
\end{align}
The $\xi$-independence of $\alpha_C$ is a consequence of the conformal
invariance of the Weyl-squared action.

\subsection{Normalized form factors at selected points}

The propagator denominators evaluate to:
\begin{equation}\label{eq:PiTT-values}
  \PiTT(0) = 1,
  \qquad
  \PiTT(1) \approx 1.179,
  \qquad
  \lim_{z \to \infty} z\,\alpha_C\,\hat{F}_1(z)
  = -\frac{89}{12}.
\end{equation}
The UV asymptotic $x\,\alpha_C\,\hat{F}_1(x) \to -89/12$ governs the
large-momentum behavior of the graviton propagator.

\subsection{Gauss--Bonnet consistency at the variational level}

The variation of the three curvature-squared invariants satisfies the
Gauss--Bonnet identity:

\begin{center}
\begin{tabular}{lcccc}
\toprule
Tensor structure & $\delta(\Riemsq)$ & $4\times\delta(\Ricsq)$
  & $\delta(R^2)$ & Sum \\
\midrule
$\Box R_{\mu\nu}$ & $4$ & $4$ & $0$ & $0$ \\
$g_{\mu\nu}\Box R$ & $4$ & $2$ & $-2$ & $0$ \\
$\nabla_\mu\nabla_\nu R$ & $-6$ & $-4$ & $2$ & $0$ \\
$R_{\mu\rho\nu\sigma}R^{\rho\sigma}$ & $8$ & $8$ & $0$ & $0$ \\
$g_{\mu\nu}\Ricsq$ & $-2$ & $-2$ & $0$ & $0$ \\
$g_{\mu\nu}R^2$ & $1/2$ & $0$ & $-1/2$ & $0$ \\
$R\,R_{\mu\nu}$ & $-2$ & $0$ & $2$ & $0$ \\
\bottomrule
\end{tabular}
\end{center}

\noindent
All seven sums vanish, confirming the Gauss--Bonnet identity
$\delta(\Riemsq) - 4\,\delta(\Ricsq) + \delta(R^2) = 0$ at the
variational level.  This identity was used to verify the dual derivation
of the field equations (Sec.~\ref{sec:field-equations}).

% ============================================================
\bibliographystyle{unsrt}
\bibliography{../references/sct_field_equations}
% ============================================================

\end{document}
